\section{Conclusion}

We introduced CoMa-20K, the first large-scale dataset pairing 3D massing geometries with functional requirements and visual urban context, enabling massing generation as a conditional VLM task. Our benchmark reveals this task's fundamental challenges: while fine-tuned models attempt complex geometries, they suffer from artifacts, whereas large zero-shot models produce clean but overly simplistic forms lacking contextual adaptation.

These findings establish CoMa-20K as a crucial benchmark and highlight two key research directions: developing specialized architectures for geometric reasoning and creating more balanced datasets. Our work provides a foundation for AI-augmented design tools that bridge programmatic requirements with contextual generation, expanding rather than replacing human creativity.